% concatenated from 05SET.tex
\documentclass[11pt]{amsart}

\usepackage[T1]{fontenc}
\usepackage[utf8]{inputenc}
\usepackage{lmodern}
\usepackage{microtype}
\usepackage{amsmath,amssymb}
\usepackage[colorlinks=true,linkcolor=blue,citecolor=blue,urlcolor=blue]{hyperref}

\newtheorem{theorem}{Theorem}[section]
\newtheorem{lemma}[theorem]{Lemma}
\newtheorem{corollary}[theorem]{Corollary}
\newtheorem{proposition}[theorem]{Proposition}
\theoremstyle{remark}
\newtheorem{remark}[theorem]{Remark}

\newcommand{\Q}{\mathbb Q}
\newcommand{\R}{\mathbb R}
\newcommand{\Z}{\mathbb Z}
\DeclareMathOperator{\ord}{ord}

\title[Neighbouring denominators and Erd\H{o}s--Mahler]
{A Neighboring-Denominator Variant of the Erd\H{o}s--Mahler Conjecture}

\author{Diego Marques}
\address{Departamento de Matem\'atica, Universidade de Bras\'ilia, Bras\'ilia, 70910-900, Brazil}
\email{diego@mat.unb.br}

\subjclass[2020]{Primary 11J70; Secondary 11J86, 11A51}
\keywords{continued fractions, Liouville numbers, smooth numbers, prime factors, $p$-adic logarithmic forms}

\begin{document}

\begin{abstract}
We prove a quantitative neighboring-denominator variant of the
Erd\H{o}s--Mahler conjecture. Let $p_n/q_n$ be the convergents of an
irrational real number $\xi$. If $p_nq_nq_{n+1}$ is $S$-smooth for
infinitely many $n$, where $S$ is a fixed finite set of primes, then
there exists an effectively computable constant $c=c(S)>0$ such that
\[
\log q_{n+1}\gg_{\xi,S} q_n^c
\]
along those indices. Consequently, $\xi$ is a Liouville number. The
proof uses the determinant identity for consecutive convergents and a
fixed-base consequence of Yu's theorem on $p$-adic logarithmic forms.
\end{abstract}


\maketitle

\section{Introduction}

Let
\[
 \xi=[a_0;a_1,a_2,\ldots]\in\R\setminus\Q
\]
and let $p_n/q_n$ denote its convergents, with $q_n>0$. For a nonzero
integer $N$, let $P(N)$ denote the largest prime divisor of $|N|$, with
$P(\pm1)=1$.

Erd\H{o}s and Mahler \cite{ErdosMahler1939} initiated the study of prime
divisors of continued-fraction convergents. In particular, they proved
that if
\begin{equation}\label{eq:EM-three}
 P(q_{n-1}q_nq_{n+1})
\end{equation}
is bounded for infinitely many $n$, then $\xi$ is a Liouville number. They further conjectured that the condition
\begin{equation}\label{eq:EM-conj}
 P(p_nq_n)\ \text{is bounded for infinitely many }n
\end{equation}
has the same consequence. We shall refer to this as the Erdős--Mahler conjecture. Related results were obtained by Fraenkel
\cite{Fraenkel1962,Fraenkel1964} and Lelis and Marques \cite{LelisMarques2017}. 

In 1983, Shorey \cite[p.~132]{Shorey1983} proved that
\[
 P(q_{n-1}p_nq_n)\longrightarrow\infty.
\]
This settles the case of the preceding denominator. The situation for the
next denominator is subtler: bounded prime support of $p_nq_nq_{n+1}$
need not itself yield a contradiction, because $q_{n+1}$ may grow very
rapidly. Our main result identifies this growth as the decisive mechanism
and shows that it already forces $\xi$ to be Liouville. 

\begin{theorem}\label{thm:main}
Let $\xi$ be an irrational real number, and let $p_n/q_n$ denote its
convergents. Suppose that $p_nq_nq_{n+1}$ is $S$-smooth for infinitely
many $n$, for some fixed finite set of primes $S$. Then there exist
$c=c(S)>0$ and $C=C(\xi,S)>0$ such that
\[
\log q_{n+1}\ge Cq_n^c
\]
for every such $n\ge2$. Consequently,
$|q_n\xi-p_n|<\exp(-Cq_n^c)$ at these indices, and in particular $\xi$
is a Liouville number.
\end{theorem}

In the terminology of \cite[Definition~2.1]{AriasKirilovMedeira2019},
Theorem~\ref{thm:main} in fact shows that $\xi$ is an exponential
Liouville number of order $1+\lceil1/c\rceil$. Its qualitative
conclusion already follows from Ridout's theorem \cite{Ridout1957}.
Indeed, if $q_{n+1}\le q_n^m$ along infinitely many relevant indices,
the determinant identity, applied to
$A_n=|p_n|q_{n+1}$ and $B_n=|p_{n+1}|q_n$, gives
$|A_n-B_n|=1$; the $S$-smoothness of $A_n$ and the polynomial bound on
the part of $B_n$ supported outside $S$ then contradict Ridout's
theorem. Thus $q_{n+1}>q_n^m$ eventually for every $m$. The new content
of Theorem~\ref{thm:main} is the exponential lower bound above.

The hypothesis is nonvacuous. There exist irrational numbers $\xi$ for
which $p_nq_nq_{n+1}$ is $3$-smooth for infinitely many $n$. To see
this, construct the continued fraction recursively. The cylinder
determined by any finite prefix contains some $2^u/3^v$, by the density
of $\{2^u/3^v:u,v\ge1\}$; extend the prefix so that
$p_n/q_n=2^u/3^v$. Since consecutive denominators are coprime,
$\gcd(q_{n-1},3)=1$, and since $2$ is a primitive root modulo $3^v$,
we may choose arbitrarily large $w$ with
$2^w\equiv q_{n-1}\pmod{3^v}$. Taking
$a_{n+1}=(2^w-q_{n-1})/3^v$ gives $q_{n+1}=2^w$, and hence
$p_nq_nq_{n+1}=2^{u+w}3^v$. Iterating gives the claim.

We now turn to the proof of Theorem~\ref{thm:main}. Its main input is a
fixed-base $p$-adic linear-form estimate, applied to the determinant
identity for consecutive convergents.

\section{Auxiliary results}

We record here the two ingredients needed for the proof of the main
theorem: standard properties of continued-fraction convergents and a
$p$-adic estimate for smooth integers.

Throughout, the notation $X\ll_{\mathcal P}Y$, equivalently
$X=O_{\mathcal P}(Y)$, means that $|X|\le CY$ for some constant
$C>0$ depending only on the parameters $\mathcal P$; the subscript is
omitted when the constant is absolute.

\subsection{Continued-fraction identities}

We begin by recalling the standard identities for continued-fraction
convergents that will be used in the proof; see, for example,
\cite[Chapter~1]{Bugeaud2004}.

\begin{equation}\label{eq:det}
 p_{n+1}q_n-p_nq_{n+1}=(-1)^n,
\end{equation}
\begin{equation}\label{eq:coprime}
 \gcd(p_n,q_n)=1,
 \qquad
 \gcd(q_n,q_{n+1})=1,
\end{equation}
together with the classical estimate
\begin{equation}\label{eq:approx}
 0<
 \left|\xi-\frac{p_n}{q_n}\right|
 <
 \frac{1}{q_nq_{n+1}}.
\end{equation}

\subsection{A $p$-adic estimate for smooth integers}

Let $S$ be a finite set of primes. A nonzero integer is said to be
\textit{$S$-smooth} if all of its prime divisors belong to $S$. For
$p$ prime and $x\in\Q^\times$, we write $\ord_p(x)$ for the normalized
$p$-adic valuation, so that $\ord_p(p)=1$.

We shall use the following fixed-base consequence of Yu's theorem \cite[\S1.1, p.~190]{Yu2007}.

\begin{lemma}\label{lem:Yu-special}
Let $p,\ell_1,\ldots,\ell_r$ be fixed primes, with $\ell_i\ne p$
for $1\le i\le r$. There exists a constant $C>0$, depending only on
$p,r,\ell_1,\ldots,\ell_r$, such that, for all integers
$b_1,\ldots,b_r$ satisfying
\[
 \ell_1^{b_1}\cdots\ell_r^{b_r}\ne1,
\]
one has
\begin{equation}\label{eq:Yu-special}
 \ord_p\!\left(
 \ell_1^{b_1}\cdots\ell_r^{b_r}-1
 \right)
 \le
 C\log\left(2+\max_{1\le i\le r}|b_i|\right).
\end{equation}
\end{lemma}

The following consequence is the form that will be applied to the
denominators of continued-fraction convergents.

\begin{lemma}\label{lem:smooth-congruence}
Let $S$ be a finite set of primes. There exists a constant $C_S>0$
such that, whenever $A,B>1$ are coprime $S$-smooth integers satisfying
\begin{equation}\label{eq:divides-abstract}
 B\mid A^2-1,
\end{equation}
one has
\begin{equation}\label{eq:smooth-congruence}
 \log B\le C_S\log\log A
\end{equation}
for all sufficiently large $A$.
\end{lemma}

\begin{proof}
Let $p\in S$ be a prime divisor of $B$. Since $\gcd(A,B)=1$, we have
$p\nmid A$, and hence
\[
 A=\prod_{\ell\in S\setminus\{p\}}\ell^{b_\ell},
 \qquad b_\ell\in\Z_{\ge0}.
\]
We apply Lemma~\ref{lem:Yu-special} to the fixed primes
$\ell\in S\setminus\{p\}$ with exponents $2b_\ell$. Since $A>1$,
\[
 \prod_{\ell\in S\setminus\{p\}}\ell^{2b_\ell}
 =A^2\ne1,
\]
so the non-vanishing hypothesis is satisfied. Therefore,
\begin{equation}\label{eq:Yu-use}
 \ord_p(A^2-1)
 \ll_{S,p}
 \log\left(
 2+\max_{\ell\in S\setminus\{p\}} b_\ell
 \right).
\end{equation}
Moreover,
\[
 b_\ell\log\ell\le\log A
\]
for every $\ell\in S\setminus\{p\}$, and hence
\[
 \max_{\ell\in S\setminus\{p\}}b_\ell
 \le \frac{\log A}{\log2}.
\]
It follows that
\begin{equation}\label{eq:vp-loglog}
 \ord_p(A^2-1)\ll_{S,p}\log\log A
\end{equation}
for all sufficiently large $A$.

Since $B\mid A^2-1$, we have
\[
 \ord_p(B)\le \ord_p(A^2-1).
\]
Summing over the primes $p\in S$ gives
\[
 \log B
 =\sum_{p\in S}\ord_p(B)\log p
 \ll_S\log\log A,
\]
which proves the lemma.
\end{proof}


We are now ready to prove the main result.


\section{Proof of Theorem~\ref{thm:main}}


Let $S$ be a finite set of primes such that $p_nq_nq_{n+1}$ is
$S$-smooth along an infinite set of indices $\mathcal N$. Since $\xi$ is irrational,
$p_n\ne0$ for all sufficiently large $n$ in this set. Put
\[
 A_n=|p_n|q_{n+1},
 \qquad
 B_n=q_n.
\]
Then $A_n$ and $B_n$ are $S$-smooth. Moreover,
\eqref{eq:coprime} gives
\[
 \gcd(A_n,B_n)=1.
\]
By \eqref{eq:det},
\[
 p_nq_{n+1}\equiv\pm1\pmod{q_n},
\]
and therefore
\begin{equation}\label{eq:central-congruence}
 q_n\mid A_n^2-1.
\end{equation}
Lemma~\ref{lem:smooth-congruence} now yields
\begin{equation}\label{eq:logqn-loglogA}
 \log q_n\ll_S\log\log A_n.
\end{equation}
Consequently, there exist constants $c,c_1>0$, depending only on $S$,
such that
\begin{equation}\label{eq:logA-power}
 \log A_n\ge c_1q_n^c
\end{equation}
for all sufficiently large $n\in\mathcal N$.

Since $p_n/q_n\to\xi$, we have $p_n=O_{\xi}(q_n)$. Hence
\[
 \log q_{n+1}
 =\log A_n-\log|p_n|
 \ge c_1q_n^c-O_\xi(\log q_n).
\]

Therefore, we obtain
\begin{equation}\label{eq:quant}
 \log q_{n+1}\ge Cq_n^c
\end{equation}
for some $C=C(\xi,S)>0$ and all sufficiently large $n\in\mathcal N$.
Consequently, for every $m\ge1$,
\[
 q_{n+1}>q_n^m
\]
whenever $n\in\mathcal N$ is sufficiently large. Hence, by
\eqref{eq:approx},
\[
 0<
 \left|\xi-\frac{p_n}{q_n}\right|
 <
 \frac{1}{q_n^{m+1}}.
\]
Since $\mathcal N$ is infinite and $m$ is arbitrary, $\xi$ is a
Liouville number.
\qed


\section*{Acknowledgements}
%The author thanks the referee for several helpful comments and suggestions that improved the presentation of the paper. 

This work was supported by the Brazilian National Council for Scientific and Technological Development (CNPq), Grant No.~304467/2023-5.


\section*{Declaration on the use of AI tools}
OpenAI's ChatGPT was used to assist
with English-language and stylistic editing of the manuscript. The author takes full responsibility for the manuscript.


\begin{thebibliography}{99}

\bibitem{AriasKirilovMedeira2019}
A.~Arias Junior, A.~Kirilov and C.~de Medeira,
Global Gevrey hypoellipticity on the torus for a class of systems of
complex vector fields,
\emph{J. Math. Anal. Appl.} \textbf{474} (2019), no.~1, 712--732.

\bibitem{Bugeaud2004}
Y.~Bugeaud,
\emph{Approximation by Algebraic Numbers},
Cambridge Tracts in Mathematics, vol.~160,
Cambridge University Press, Cambridge, 2004.

\bibitem{ErdosMahler1939}
P.~Erd\H{o}s and K.~Mahler,
Some arithmetical properties of the convergents of a continued fraction,
\emph{J. London Math. Soc.} \textbf{14} (1939), 12--18.

\bibitem{Fraenkel1962}
A.~S. Fraenkel,
On a theorem of D. Ridout in the theory of Diophantine approximations,
\emph{Trans. Amer. Math. Soc.} \textbf{105} (1962), 84--101.

\bibitem{Fraenkel1964}
A.~S. Fraenkel,
Transcendental numbers and a conjecture of Erd\H{o}s and Mahler,
\emph{J. London Math. Soc.} \textbf{39} (1964), 405--416.

\bibitem{LelisMarques2017}
J.~Lelis and D.~Marques,
On a problem of Erd\H{o}s and Mahler concerning continued fractions,
\emph{Bull. Aust. Math. Soc.} \textbf{95} (2017), no.~2, 183--186.

\bibitem{Ridout1957} D. Ridout, Rational approximations to algebraic numbers, \textit{Mathematika} \textbf{4} (1957), 125--131.

\bibitem{Shorey1983}
T.~N. Shorey,
Divisors of convergents of a continued fraction,
\emph{J. Number Theory} \textbf{17} (1983), 127--133.

\bibitem{Yu2007}
K.~Yu,
$p$-adic logarithmic forms and group varieties. III,
\emph{Forum Math.} \textbf{19} (2007), no.~2, 187--280.

\end{thebibliography}



\end{document}
